\documentclass[11pt,letterpaper]{article}
\usepackage[T1]{fontenc}
\usepackage[utf8]{inputenc}
\usepackage{amsmath,amsthm,mathtools}
\usepackage{newtxtext,newtxmath}
\usepackage[scaled=0.92]{sourcesanspro}
\usepackage{microtype}
\usepackage[margin=1in,headheight=15pt]{geometry}
\usepackage{booktabs,array,longtable}
\usepackage{xcolor}
\usepackage[most]{tcolorbox}
\usepackage{enumitem}
\usepackage{listings}
\usepackage{fancyhdr}
\usepackage{titlesec}
\usepackage{hyperref}
\usepackage[nameinlink,noabbrev]{cleveref}
\definecolor{Olive}{HTML}{526347}
\definecolor{PaleOlive}{HTML}{F1F3EA}
\definecolor{Ink}{HTML}{263126}
\definecolor{Muted}{HTML}{5E675D}
\definecolor{Rule}{HTML}{C6CDBD}
\hypersetup{colorlinks=true,linkcolor=Olive,citecolor=Olive,urlcolor=Olive,
  pdftitle={Rationality and growth of stabilized permutation-weight series},
  pdfauthor={OpenAI ChatGPT},pdfsubject={Exact generating functions and recurrences}}
\titleformat{\section}{\Large\sffamily\bfseries\color{Olive}}{\thesection}{0.65em}{}
\titleformat{\subsection}{\large\sffamily\bfseries\color{Ink}}{\thesubsection}{0.65em}{}
\titlespacing*{\section}{0pt}{2.2ex plus .7ex}{1ex}
\titlespacing*{\subsection}{0pt}{1.8ex plus .5ex}{.8ex}
\pagestyle{fancy}
\fancyhf{}
\fancyhead[L]{\small\sffamily\color{Muted}Stabilized permutation-weight series}
\fancyhead[R]{\small\sffamily\color{Muted}September 2026}
\fancyfoot[C]{\small\sffamily\thepage}
\renewcommand{\headrulewidth}{0.3pt}
\setlength{\parindent}{1.25em}
\setlength{\parskip}{3pt}
\setlength{\emergencystretch}{2em}
\allowdisplaybreaks[2]
\numberwithin{equation}{section}
\newtheorem{theorem}{Theorem}[section]
\newtheorem{proposition}[theorem]{Proposition}
\newtheorem{lemma}[theorem]{Lemma}
\newtheorem{corollary}[theorem]{Corollary}
\theoremstyle{definition}
\newtheorem{definition}[theorem]{Definition}
\newtheorem{problem}[theorem]{Problem}
\theoremstyle{remark}
\newtheorem{remark}[theorem]{Remark}
\newcommand{\Q}{\mathbb Q}
\newcommand{\Z}{\mathbb Z}
\newcommand{\N}{\mathbb N}
\newcommand{\A}{A}
\newcommand{\Spec}{\Lambda}
\newcommand{\coef}[2]{[#1]#2}
\newcommand{\ord}{\operatorname{ord}}
\newcommand{\res}{\operatorname{Res}}
\newcommand{\des}{\operatorname{des}}
\newcommand{\wt}{\operatorname{wt}}
\newtcolorbox{resultbox}[1][]{enhanced,breakable,colback=PaleOlive,colframe=Rule,
  boxrule=.5pt,arc=1.3mm,left=3mm,right=3mm,top=2mm,bottom=2mm,#1}
\lstset{basicstyle=\small\ttfamily,columns=fullflexible,keepspaces=true,
  frame=single,rulecolor=\color{Rule},backgroundcolor=\color{PaleOlive},
  breaklines=true,showstringspaces=false,aboveskip=10pt,belowskip=10pt}

\begin{document}
\thispagestyle{empty}
\begin{center}
  {\small\sffamily\bfseries\color{Olive}EXACT COMBINATORICS \quad / \quad GENERATING FUNCTIONS}\par
  \vspace{12pt}
  {\LARGE\bfseries Rationality and growth of stabilized\\[3pt]
  permutation-weight series}\par
  \vspace{9pt}
  {\large Finite spectra, integer recurrences, and an explicit two-descent formula}\par
  \vspace{12pt}
  {\sffamily Research report prepared for Vladimir Reshetnikov}\par
  \vspace{4pt}
  {\small September 19, 2026}
\end{center}
\vspace{8pt}
\begin{abstract}
The permutation-weight $q$-Eulerian polynomials have stabilized leading
coefficients $a_{d,k}$ when the number of descents $d$ is fixed. Earlier work
proved stabilization and investigated an initial agreement with the partition
triangle OEIS A256193. We pursue the complete-series problem: can the entire
series $W_d(t)=\sum_{k\ge0}a_{d,k}t^k$, including the coefficients beyond that
agreement, be described by a finite exact formula?

Starting from the published polynomial recurrence, we prove that every $W_d$
is rational. The key is a finite exponential spectrum indexed by partitions
of integers at most $d+1$. We give an explicit universal denominator, an
exact construction of the numerator, and a practical symbolic algorithm.
We also prove that the radius of convergence is exactly
$(d+1)^{-1/d}$, with a periodic leading asymptotic term. For $d=2$ we obtain
closed formulas involving only $3^m$, $2^m$, and Fibonacci numbers.
In contrast, the fixed-$d$ diagonal series of A256193 is nonrational and
has radius~$1$. Exact computations for $1\le d\le6$ are independently checked
against the original recurrence through deficit~$40$; the archive contains
all code, certificates, rational functions, and coefficient data.
\end{abstract}

\begin{resultbox}
\small
\textbf{Scope and status.}
The stabilization theorem itself is known, and the earlier near-diagonal
partition correspondence is not claimed as new. The problem selected here
is a precise generating-function continuation of the request for further
coefficients in Huang--Tang: determine each \emph{whole} stabilized series by
finite algebraic data. The rationality theorem, denominator construction,
and growth analysis below were derived in this investigation and were not
found in the three primary papers reviewed. This is a proved research
result within the report, not an exhaustive certification of novelty or a
claim of peer review. The English proofs do not depend on numerical fitting
or on an unproved conjecture.
\end{resultbox}

\vspace{3pt}
\begin{center}
\small\sffamily
Companion files: \texttt{code/derive.py}, \texttt{code/verify.py},
\texttt{data/rational\_series.json}.
\end{center}

\newpage
\begingroup
\small
\setlength{\parskip}{0pt}
\tableofcontents
\endgroup
\vspace{4pt}
\begin{resultbox}[breakable=false]
\textbf{The principal conclusions at a glance.}
For every fixed integer $d\ge1$,
\[
 W_d(t)\in\Q(t)\cap\Z_{\ge0}[[t]],\qquad
 \operatorname{rad}(W_d)=(d+1)^{-1/d},\qquad
 \limsup_{k\to\infty}a_{d,k}^{1/k}=(d+1)^{1/d}.
\]
For the two-descent sequence, with $F_0=0$, $F_1=1$,
\begin{align*}
 a_{2,2m}&=45\,3^m+8\,2^m-F_{2m+10}+m+3,\\
 a_{2,2m+1}&=81\,3^m+8\,2^m-F_{2m+11}+m+4.
\end{align*}
These identities hold for every $m\ge0$, not merely over a tested range.
\end{resultbox}
\clearpage

\section{The selected problem and its provenance}
\subsection{A specific family, not an arbitrary \texorpdfstring{$q$}{q}-Eulerian convention}
There are several inequivalent $q$-analogues of Eulerian polynomials.
Here the family is the permutation-weight family introduced by
Dugan, Glennon, Gunnells, and Steingr\'imsson \cite{DGGs}.
Its terms enumerate permutations by descents and a recursively defined
weight arising from maxmin trees. To remove any ambiguity about the
statistic or normalization, we take the recurrence established by
Agrawal, Choi, and Sun \cite[Theorem 4.1]{ACS} as the exact algebraic
specification:
\begin{equation}\label{eq:original}
\begin{split}
 E_0(x,q)&=1,\\
 E_n(x,q)&=E_{n-1}(qx,q)\\
 &\quad+x\sum_{i=1}^{n-1}\binom{n-1}{i}
              E_i(x,q)E_{n-i-1}(qx,q)\qquad(n\ge1).
\end{split}
\end{equation}
Write
\begin{equation}\label{eq:Adefinition}
 A_{n,d}(q)=[x^d]E_n(x,q),\qquad D(n,d)=d(n-d-1).
\end{equation}
For $n\ge d+1$, the degree of $A_{n,d}$ is $D(n,d)$ and its
leading coefficient is~$1$. Reversing the polynomial gives
\begin{equation}\label{eq:Rdefinition}
 R_{n,d}(t)=t^{D(n,d)}A_{n,d}(t^{-1}).
\end{equation}
The coefficient of $t^k$ counts weight deficit $k$ below the largest weight
available at length $n$ and descent count~$d$.

The coefficients stabilize as $n$ increases. In fact, the coefficient of
$t^k$ is fixed from $n=d+k+1$ onward. Define
\begin{equation}\label{eq:Wdefinition}
 a_{d,k}=[t^k]R_{d+k+1,d}(t),\qquad
 W_d(t)=\sum_{k\ge0}a_{d,k}t^k\qquad(d\ge1).
\end{equation}
This stabilization is due to Agrawal--Choi--Sun \cite{ACS}.
We reproduce a short proof in \cref{sec:stabilization}, both to make the
argument self-contained from \eqref{eq:original} and to justify the
independent verification algorithm.

\subsection{What was already known, and what is investigated here}
Dugan et al. observed stabilization and a partial agreement with the
partition triangle A256193 \cite[Section 6.12]{DGGs}. Agrawal et al. proved
stabilization and proposed a recurrence relating certain near-diagonal
coefficients \cite[Section 6]{ACS}. Huang and Tang subsequently gave a
maxmin-tree treatment of the partition connection; their final discussion
asks for ways to reach further coefficients when more complicated tree
structures enter \cite[Section 5]{HT}. Thus it would be incorrect to
present stabilization or the original partition correspondence as an
unresolved conjecture in this report.

The question pursued here is a more specific continuation:
\begin{problem}\label{prob:whole}
For fixed $d$, determine whether the \emph{entire} stabilized sequence
$(a_{d,k})_{k\ge0}$ satisfies a finite constant-coefficient linear
recurrence. If it does, construct one exactly, and determine its
exponential growth.
\end{problem}
This formulation is ours; it is not attributed to the earlier authors
as a verbatim named conjecture. It asks for information substantially
stronger than computing a finite prefix from \eqref{eq:original}.
The answer will be affirmative for every~$d$.

\section{Leading terms and stabilization}\label{sec:stabilization}
Taking the coefficient of $x^d$ in \eqref{eq:original} gives
\begin{equation}\label{eq:Arecurrence}
 A_{n,d}=q^d A_{n-1,d}
 +\sum_{i=1}^{n-1}\binom{n-1}{i}
       \sum_{a+b=d-1}q^b A_{i,a}A_{n-i-1,b}.
\end{equation}
All omitted arguments are $q$. We use $A_{0,0}=1$, $A_{0,d}=0$ for $d>0$,
and $A_{n,d}=0$ outside $0\le d\le n-1$ when $n\ge1$.

\begin{lemma}\label{lem:degree}
For $n\ge d+1$ and $d\ge1$, $A_{n,d}$ has nonnegative integer
coefficients, degree $D(n,d)$, and leading coefficient~$1$.
Also $A_{n,0}=1$ for every $n\ge0$.
\end{lemma}
\begin{proof}
Positivity and the assertion for $d=0$ follow immediately by induction.
The support bound $d\le n-1$ also follows from \eqref{eq:Arecurrence}.
At $n=d+1$, the only nonzero convolution term is
$i=d$, $a=d-1$, and an empty right factor, so $A_{d+1,d}=1$.

For the degree comparison, let $m=n-i-1$. First suppose $m>0$.
For a nonzero term write
\[
 i=a+1+u,\quad m=b+1+v,\quad u,v\ge0,\quad a+b=d-1.
\]
The difference between $D(n,d)$ and the maximal degree of
$q^bA_{i,a}A_{m,b}$ is
\begin{equation}\label{eq:gap}
 \Delta=a+1+(b+1)u+(a+1)v.
\end{equation}
This is positive. If $m=0$, then $b=0$, $a=d-1$, and the same difference is
\begin{equation}\label{eq:emptygap}
 \Delta=n-d-1.
\end{equation}
For $n>d+1$ this is again positive. On the other hand,
$q^dA_{n-1,d}$ has degree exactly $D(n,d)$ and leading coefficient~$1$.
It is therefore the unique contribution to the leading term.
\end{proof}

\begin{proposition}[Known stabilization, reproved]\label{prop:stabilization}
For $n\ge d+2$,
\begin{equation}\label{eq:stabilization}
 R_{n,d}(t)-R_{n-1,d}(t)\in t^{n-d-1}\Z_{\ge0}[t].
\end{equation}
Consequently $[t^k]R_{n,d}$ is independent of $n$ for $n\ge d+k+1$.
\end{proposition}
\begin{proof}
After reversal, the first term of \eqref{eq:Arecurrence} is exactly
$R_{n-1,d}$. A nonempty-right convolution term becomes a product of
reversed polynomials multiplied by $t^\Delta$, where \eqref{eq:gap} gives
\[
 \Delta-(n-d-1)
 =\bigl(a+1+(b+1)u+(a+1)v\bigr)-(1+u+v)
 =a+bu+av\ge0.
\]
An empty-right term has exactly the shift in \eqref{eq:emptygap}.
All coefficients are nonnegative. Thus the increment at length $n$
vanishes in degrees below $n-d-1$. Once $n=d+k+1$ has been reached,
every subsequent increment starts in degree at least $k+1$.
\end{proof}

The limit $R_{n,d}\to W_d$ is a \emph{formal} power-series limit:
each fixed coefficient eventually stops changing. No numerical value of
$t$, analytic convergence, or interchange of infinite analytic sums is
needed at this stage.

\section{The finite exponential spectrum}\label{sec:spectrum}
\subsection{The differential hierarchy}
Introduce exponential generating functions in the length variable:
\begin{equation}\label{eq:egf}
 \mathcal F(z,x,q)=\sum_{n\ge0}E_n(x,q)\frac{z^n}{n!},\qquad
 f_d(z,q)=[x^d]\mathcal F(z,x,q).
\end{equation}
All calculations in this section are over the field $K=\Q(q)$.
The binomial convolution in \eqref{eq:original} becomes multiplication:
\begin{equation}\label{eq:functionalegf}
 \partial_z\mathcal F(z,x,q)
 =\bigl(1+x(\mathcal F(z,x,q)-1)\bigr)\mathcal F(z,qx,q).
\end{equation}
Hence $f_0=e^z$, while, for $d\ge1$,
\begin{equation}\label{eq:ode}
 \partial_z f_d=q^d f_d+g_d,\qquad f_d(0,q)=0,
\end{equation}
where
\begin{equation}\label{eq:g}
 g_d(z,q)=\sum_{a+b=d-1}q^b f_a(z,q)f_b(z,q)
                    -q^{d-1}f_{d-1}(z,q).
\end{equation}
The subtracted term is essential: it removes the forbidden empty left
factor in the original recurrence. Omitting it produces a different
sequence.

\subsection{A partition-indexed set of possible frequencies}
\begin{definition}\label{def:spectrum}
For $d\ge0$, let
\begin{equation}\label{eq:spectrum}
 \Spec_d=\left\{
 \lambda_{\boldsymbol c}(q)=\sum_{j=0}^{d}c_jq^j:
 c_j\in\Z_{\ge0},\quad 1\le\sum_{j=0}^{d}(j+1)c_j\le d+1
 \right\}.
\end{equation}
The integer $\sum(j+1)c_j$ is the cost of the frequency.
Write $M_d=|\Spec_d|$.
\end{definition}
A vector $\boldsymbol c$ of cost $r$ is exactly a partition of $r$ with
$c_j$ parts of size $j+1$. Distinct vectors give distinct polynomials in
$q$. Therefore
\begin{equation}\label{eq:count}
 M_d=\sum_{r=1}^{d+1}p(r),
\end{equation}
where $p(r)$ is the ordinary partition number. The first values are
$M_0,M_1,\ldots,M_6=1,3,6,11,18,29,44$.

For example,
\begin{align*}
 \Spec_1&=\{1,2,q\},\\
 \Spec_2&=\{1,2,3,q,q+1,q^2\},\\
 \Spec_3&=\{1,2,3,4,q,q+1,q+2,2q,q^2,q^2+1,q^3\}.
\end{align*}
The only frequency in $\Spec_d$ of degree $d$ is $q^d$.
Every other frequency has degree at most $d-1$.

\begin{theorem}[Finite-spectrum theorem]\label{thm:spectrum}
For every $d\ge0$, there are rational functions
$c_{d,\lambda}(q)\in\Q(q)$ such that
\begin{equation}\label{eq:expansion}
 f_d(z,q)=\sum_{\lambda\in\Spec_d}c_{d,\lambda}(q)e^{\lambda(q)z}.
\end{equation}
In particular,
\begin{equation}\label{eq:powers}
 A_{n,d}(q)=\sum_{\lambda\in\Spec_d}c_{d,\lambda}(q)\lambda(q)^n
 \qquad(n\ge0).
\end{equation}
Some allowed coefficients may be zero; the theorem does not assert
that every allowed frequency occurs for every~$d$.
\end{theorem}
\begin{proof}
For $d=0$, $f_0=e^z$ and $\Spec_0=\{1\}$.
Suppose the assertion holds at every smaller descent count.
If $a+b=d-1$, a frequency from $f_af_b$ is a sum
$\lambda+\mu$ of costs at most $(a+1)+(b+1)=d+1$.
It belongs to $\Spec_d$. Frequencies from $f_{d-1}$ also belong to
$\Spec_d$. None of these has degree $d$: both factors in a product
have descent count at most $d-1$.

Write $g_d=\sum_\lambda h_{d,\lambda}e^{\lambda z}$ after collecting
equal frequencies. Solving \eqref{eq:ode} gives the finite expression
\begin{equation}\label{eq:solveode}
 f_d(z,q)=\sum_\lambda
       \frac{h_{d,\lambda}(q)}{\lambda(q)-q^d}
       \bigl(e^{\lambda(q)z}-e^{q^dz}\bigr).
\end{equation}
Every denominator is nonzero in $\Q(q)$, because $\deg\lambda<d$.
The only new frequency is $q^d$, whose cost is $d+1$.
This proves \eqref{eq:expansion}; coefficient extraction in $z$ gives
\eqref{eq:powers}.
\end{proof}

Distinct exponentials over $K$ are linearly independent: comparing the
first $M_d$ derivatives at $z=0$ gives a Vandermonde matrix with nonzero
determinant. Thus the collected coefficients in \eqref{eq:expansion}
are unique. The absence of resonance in \eqref{eq:solveode} is an identity
in the indeterminate $q$, not an assertion that particular numerical
specializations of $q$ can never produce coincident frequencies.

\section{Rationality and a universal integer denominator}\label{sec:rationality}
Let
\begin{equation}\label{eq:C}
 C_d(q)=c_{d,q^d}(q).
\end{equation}
The coefficient $C_d$ is distinguished by degree, not by ordering
frequencies at a particular real value of~$q$.

\begin{theorem}[Rationality of the complete stabilized series]\label{thm:rationality}
For every integer $d\ge1$,
\begin{equation}\label{eq:mainW}
 \boxed{\displaystyle W_d(t)=t^{-d(d+1)}C_d(t^{-1}).}
\end{equation}
Consequently $W_d\in\Q(t)\cap\Z_{\ge0}[[t]]$.
\end{theorem}
\begin{proof}
Substitute \eqref{eq:powers} into \eqref{eq:Rdefinition}:
\begin{equation}\label{eq:limit}
 R_{n,d}(t)=t^{-d(d+1)}
 \sum_{\lambda\in\Spec_d}c_{d,\lambda}(t^{-1})
                  \bigl(t^d\lambda(t^{-1})\bigr)^n.
\end{equation}
For $\lambda=q^d$, the parenthesized factor is~$1$.
For every other frequency it is a polynomial divisible by $t$,
since $\deg\lambda\le d-1$.
Each rational amplitude, viewed in $\Q((t))$, has a finite Laurent
principal part. Multiplication by the $n$th power of a polynomial of
positive $t$-adic order therefore tends coefficientwise to zero.
There are finitely many frequencies, so only the $q^d$ term survives.
By \cref{prop:stabilization}, the limit on the left is $W_d$.
Its integrality, nonnegativity, and regularity at $t=0$ were already
established directly from the recurrence.
\end{proof}

\subsection{A denominator requiring no guesswork}
Define
\begin{equation}\label{eq:univden}
 \mathcal D_d(t)=
 \prod_{\lambda\in\Spec_d\setminus\{q^d\}}
       \left(1-t^d\lambda(t^{-1})\right).
\end{equation}
Every factor is an integer polynomial with constant coefficient~$1$.

\begin{theorem}[Explicit denominator and finite numerator]\label{thm:denominator}
For every $d\ge1$, there is a polynomial $\mathcal P_d(t)\in\Z[t]$ with
\begin{equation}\label{eq:univrepr}
 W_d(t)=\frac{\mathcal P_d(t)}{\mathcal D_d(t)}.
\end{equation}
The numerator can be constructed from the finitely many polynomials
$A_{0,d},\ldots,A_{M_d-1,d}$ as follows. Put
\begin{equation}\label{eq:annihilator}
 B_d(u,q)=\prod_{\lambda\in\Spec_d\setminus\{q^d\}}(u-\lambda(q))
            =\sum_{j=0}^{M_d-1}b_{d,j}(q)u^j,
\end{equation}
then set
\begin{equation}\label{eq:numeratorformula}
\begin{split}
 N_d(q)&=\sum_{j=0}^{M_d-1}b_{d,j}(q)A_{j,d}(q),\\
 \mathcal P_d(t)&=t^{d(M_d-1)-d(d+1)}N_d(t^{-1}).
\end{split}
\end{equation}
Although the last expression is initially written as a Laurent
polynomial, all its negative powers cancel.
\end{theorem}
\begin{proof}
Apply the polynomial $B_d$ to the power-sum representation
\eqref{eq:powers} at $n=0$. It annihilates every frequency except $q^d$:
\[
 N_d(q)=C_d(q)B_d(q^d,q).
\]
Thus
\begin{equation}\label{eq:Cformula}
 C_d(q)=\frac{N_d(q)}{
          \prod_{\lambda\ne q^d}(q^d-\lambda(q))}.
\end{equation}
Substitution into \eqref{eq:mainW} gives \eqref{eq:univrepr} and
\eqref{eq:numeratorformula}. All $b_{d,j}$ and $A_{j,d}$ are integer
polynomials, so the expression for $\mathcal P_d$ is a Laurent
polynomial with integer coefficients. But
$\mathcal P_d=\mathcal D_dW_d$, with $\mathcal D_d\in\Z[t]$ and
$W_d\in\Z[[t]]$. It has no negative powers, proving the claim.
\end{proof}

\begin{corollary}[An explicit constant-coefficient recurrence]\label{cor:recurrence}
Write $\mathcal D_d(t)=1+\delta_1t+\cdots+\delta_Lt^L$.
Then, for every $k>\deg\mathcal P_d$,
\begin{equation}\label{eq:coeffrecurrence}
 a_{d,k}+\sum_{j=1}^{L}\delta_j a_{d,k-j}=0,
\end{equation}
where coefficients with negative indices are set equal to zero.
In particular, a conventional recurrence in the preceding $L$ terms
holds once $k\ge\max(L,\deg\mathcal P_d+1)$.
\end{corollary}
\begin{proof}
Compare coefficients in $\mathcal D_dW_d=\mathcal P_d$.
\end{proof}

This is a constructive answer to \cref{prob:whole}. The order bound
$L\le d(M_d-1)$ is finite and explicit. It need not be minimal: cancellation
already occurs at $d=4$. For computation, the sparse differential
construction \eqref{eq:solveode} is considerably more economical than
forming every $A_{j,d}$ required by \eqref{eq:numeratorformula}.
The latter formula is especially useful as a proof of the denominator
bound and as a conceptually independent construction.

\section{Explicit results for small descent counts}\label{sec:examples}
\subsection{One descent}
At $d=1$, the differential equation is
$f_1'=qf_1+e^{2z}-e^z$. Hence
\[
 f_1(z,q)=\frac{e^{2z}-e^{qz}}{2-q}
          -\frac{e^z-e^{qz}}{1-q},\qquad
 C_1(q)=\frac{1}{(q-1)(q-2)}.
\]
Therefore
\begin{equation}\label{eq:W1}
 W_1(t)=\frac{1}{(1-t)(1-2t)},\qquad a_{1,k}=2^{k+1}-1.
\end{equation}
This elementary case also checks the reversal exponent in
\eqref{eq:mainW}.

\subsection{Two descents: the complete sequence in closed form}
The next forcing term is
\[
 g_2(z,q)=\bigl((q+1)e^z-q\bigr)f_1(z,q).
\]
For a fully explicit extraction, write
\[
 f_1=\frac{e^z}{q-1}-\frac{e^{2z}}{q-2}
           +\frac{e^{qz}}{(q-1)(q-2)}.
\]
The five forcing frequencies are $1,2,3,q,q+1$. Their amplitudes are,
in that order,
\[
 -\frac q{q-1},\quad
 \frac{q+1}{q-1}+\frac q{q-2},\quad
 -\frac{q+1}{q-2},\quad
 -\frac q{(q-1)(q-2)},\quad
 \frac{q+1}{(q-1)(q-2)}.
\]
Summing each amplitude divided by $q^2-\lambda$, as prescribed by
\eqref{eq:solveode}, and simplifying gives
\[
 C_2(q)=\frac{q^3+2q^2-3q-2}
 {(q-1)^2(q+1)(q^2-2)(q^2-3)(q^2-q-1)}.
\]
Consequently
\begin{equation}\label{eq:W2}
 \boxed{\displaystyle
 W_2(t)=\frac{1+2t-3t^2-2t^3}
 {(1-t)^2(1+t)(1-2t^2)(1-3t^2)(1-t-t^2)}.}
\end{equation}
No interpolation of initial sequence values is used to obtain this
expression: it is the exact result of solving the two differential
levels.

A useful partial-fraction identity is
\begin{equation}\label{eq:W2apart}
\begin{split}
 W_2(t)={}&\frac{8(1+t)}{1-2t^2}
       +\frac{9(5+9t)}{1-3t^2}
       -\frac{55+34t}{1-t-t^2}\\
       &-\frac{1}{4(1+t)}+\frac{11}{4(1-t)}
        +\frac{1}{2(1-t)^2}.
\end{split}
\end{equation}
It can be checked simply by multiplication by the denominator in
\eqref{eq:W2}.

\begin{theorem}[Exact two-descent formula]\label{thm:W2closed}
Let $F_0=0$, $F_1=1$, and $F_{n+2}=F_{n+1}+F_n$. For every $m\ge0$,
\begin{align}
 a_{2,2m}&=45\,3^m+8\,2^m-F_{2m+10}+m+3,\label{eq:even}\\
 a_{2,2m+1}&=81\,3^m+8\,2^m-F_{2m+11}+m+4.\label{eq:odd}
\end{align}
\end{theorem}
\begin{proof}
The first two terms of \eqref{eq:W2apart} contribute $8\,2^m$ in both
parities, and respectively $45\,3^m$ and $81\,3^m$.
The coefficient of the third term is
$-55F_{k+1}-34F_k=-F_{k+10}$, using the Fibonacci addition identity.
The last three terms contribute
\[
 -\frac{(-1)^k}{4}+\frac{11}{4}+\frac{k+1}{2},
\]
which is $m+3$ for $k=2m$ and $m+4$ for $k=2m+1$.
\end{proof}

The expanded denominator is
\[
 Q_2(t)=1-2t-6t^2+13t^3+11t^4-28t^5-6t^6+23t^7-6t^9.
\]
Thus the ordinary recurrence of order nine is
\begin{equation}\label{eq:W2recurrence}
\begin{split}
 a_{2,k}={}&2a_{2,k-1}+6a_{2,k-2}-13a_{2,k-3}-11a_{2,k-4}\\
 &+28a_{2,k-5}+6a_{2,k-6}-23a_{2,k-7}+6a_{2,k-9},
 \qquad k\ge9.
\end{split}
\end{equation}
Its initial values are
\[
 (a_{2,0},\ldots,a_{2,8})=(1,4,11,31,65,157,298,661,1196).
\]
The order is minimal. Indeed, the numerator in \eqref{eq:W2} is nonzero
at $t=1,-1$; modulo $1-2t^2$, $1-3t^2$, and $1-t-t^2$ its remainders are,
respectively, $t-\tfrac12$, $\tfrac43t$, and $t$.
None shares a zero with its corresponding quadratic. Therefore numerator
and denominator are coprime, and no smaller eventual constant-coefficient
recurrence can have this rational generating function.

An immediate asymptotic consequence is
\begin{equation}\label{eq:W2asymptotic}
 a_{2,2m}\sim45\,3^m,\qquad a_{2,2m+1}\sim81\,3^m.
\end{equation}
The Fibonacci error is smaller because $\varphi^2=(3+\sqrt5)/2<3$.
In particular, $a_{2,k}^{1/k}\to\sqrt3$, but consecutive-term ratios do
\emph{not} converge:
\[
 \frac{a_{2,2m+1}}{a_{2,2m}}\longrightarrow\frac95,
 \qquad
 \frac{a_{2,2m+2}}{a_{2,2m+1}}\longrightarrow\frac53.
\]
A single-ratio numerical extrapolation would obscure this parity effect.

\subsection{Three descents and beyond}
For $d=3$ no cancellation of the universal denominator occurs. One obtains
$W_3=P_3/Q_3$, where
\begin{equation}\label{eq:Q3}
\begin{split}
 Q_3(t)={}&(1-t)(1-t^2)(1-2t^2)
 (1-t^3)(1-2t^3)(1-3t^3)(1-4t^3)\\
 &\quad\cdot(1-t^2-t^3)(1-t^2-2t^3)(1-t-t^3),
\end{split}
\end{equation}
\begin{equation}\label{eq:P3}
\begin{split}
 P_3(t)={}&1+3t+2t^2-15t^3-23t^4-17t^5+71t^6+74t^7\\
 &+23t^8-117t^9-86t^{10}-26t^{11}+54t^{12}+44t^{13}+24t^{14}.
\end{split}
\end{equation}
The resulting recurrence has order~$26$.
Large expanded polynomials become less readable at higher $d$, so their
complete coefficient arrays are supplied in the companion JSON rather
than typeset as several pages of integers.

\begin{table}[htbp]
\centering
\caption{Exact symbolic results. $Q_d$ is the reduced denominator with
$Q_d(0)=1$; $P_d$ is its coprime numerator.}
\label{tab:degrees}
\begin{tabular}{rrrrr}
\toprule
$d$ & $M_d$ & $\deg\mathcal D_d$ & $\deg Q_d$ & $\deg P_d$\\
\midrule
1 & 3  & 2   & 2   & 0\\
2 & 6  & 9   & 9   & 3\\
3 & 11 & 26  & 26  & 14\\
4 & 18 & 60  & 57  & 37\\
5 & 29 & 124 & 121 & 91\\
6 & 44 & 233 & 223 & 181\\
\bottomrule
\end{tabular}
\end{table}
The exact canceled factors for the last three rows are
\begin{align*}
 \mathcal D_4/Q_4=\mathcal D_5/Q_5&=(1+t)(1+t^2),\\
 \mathcal D_6/Q_6&=(1+t)^4(1+t^2)(1-t+t^2)(1+t+t^2).
\end{align*}
These cancellations are why the universal denominator should not be
advertised as the minimal denominator in general.

\section{The exact radius of convergence}\label{sec:growth}
Rationality gives a finite set of possible poles. We next locate the
nearest ones and prove that the positive real candidate does not cancel.
Put
\begin{equation}\label{eq:rho}
 r_d=(d+1)^{1/d},\qquad \rho_d=r_d^{-1}.
\end{equation}

\begin{lemma}[Spectral separation at the first candidate pole]\label{lem:separation}
For $\lambda\in\Spec_d\setminus\{q^d\}$,
\[
 \lambda(r_d)\le d+1.
\]
Equality holds only for the constant frequency $\lambda=d+1$.
Consequently every factor of $\mathcal D_d(t)$ is nonzero for
$|t|<\rho_d$. On $|t|=\rho_d$, the only possible zeros arise from
$1-(d+1)t^d$.
\end{lemma}
\begin{proof}
The function $\log(1+x)/x$ is strictly decreasing for $x>0$.
Indeed, its derivative has the sign of
$x/(1+x)-\log(1+x)<0$; the latter inequality follows by integrating
$1/(1+u)>1/(1+x)$ on $0<u<x$.
Thus, for $1\le j<d$,
\[
 r_d^j=(d+1)^{j/d}<j+1.
\]
For $j=0$, equality $r_d^0=1$ holds.
A frequency other than $q^d$ has $c_d=0$, hence
\[
 \lambda(r_d)=\sum_{j<d}c_jr_d^j
       \le\sum_{j<d}(j+1)c_j\le d+1.
\]
Equality requires that only $c_0$ be nonzero and $c_0=d+1$.

For complex $t$ with $|t|\le\rho_d$, the triangle inequality gives
\[
 |t^d\lambda(t^{-1})|
 \le\sum_{j<d}c_j|t|^{d-j}
 \le r_d^{-d}\lambda(r_d)\le1.
\]
The inequality is strict in the open disk, and also on its boundary
unless $\lambda=d+1$. This proves the pole assertion.
\end{proof}

\subsection{Why the positive pole survives}
Let $h_j(q)$ be the coefficient of $e^{(j+1)z}$ in $f_j(z,q)$.
The finite-spectrum construction gives $h_0=1$ and
\begin{equation}\label{eq:hrecurrence}
 h_j(q)=\frac{H_j(q)}{j+1-q^j},\qquad
 H_j(q)=\sum_{a+b=j-1}q^b h_a(q)h_b(q)\quad(j\ge1).
\end{equation}
To see this, a product $f_af_b$ has constant frequency $j+1$ only when
its factors use their largest constant frequencies $a+1$ and $b+1$.
The subtracted term in \eqref{eq:g} has no such frequency.

\begin{lemma}[Noncancellation]\label{lem:noncancel}
$C_d(q)$ has a simple pole at $q=r_d$, with residue
\begin{equation}\label{eq:residue}
 \res_{q=r_d}C_d(q)=\frac{H_d(r_d)}{d\,r_d^{d-1}}>0.
\end{equation}
Equivalently, the rational function
\begin{equation}\label{eq:L}
 L_d(t)=\bigl(1-(d+1)t^d\bigr)W_d(t)
\end{equation}
is regular at $t=\rho_d$ and
\begin{equation}\label{eq:Lpositive}
 L_d(\rho_d)=r_d^{d^2}H_d(r_d)>0.
\end{equation}
\end{lemma}
\begin{proof}
For $j<d$, the denominator $j+1-r_d^j$ in \eqref{eq:hrecurrence}
is positive. Induction gives $h_j(r_d)>0$, so $H_d(r_d)>0$.
It follows that $h_d(q)$ has a simple pole at $q=r_d$, of residue
$-H_d(r_d)/(d r_d^{d-1})$.

On the other hand, \eqref{eq:Cformula} and
\cref{lem:separation} show that $C_d$ has at most a simple pole there:
only the denominator factor $q^d-(d+1)$ vanishes, and its derivative is
$d r_d^{d-1}\ne0$.

The singular parts of the terms
$C_d(q)e^{q^dz}$ and $h_d(q)e^{(d+1)z}$ must cancel.
Here is a justification that also allows collisions among other
frequencies. Every Taylor coefficient of $f_d$ in $z$ is a polynomial
in $q$, so the coefficient of $(q-r_d)^{-1}$ in its finite exponential
representation is identically zero. For any other frequency, that
Laurent coefficient is a polynomial in $z$ times
$e^{\lambda(r_d)z}$, obtained by finitely many Taylor terms of the
exponential. All these exponents satisfy $\lambda(r_d)<d+1$.
Polynomials times exponentials at distinct exponents are linearly
independent: applying sufficiently high powers of
$\partial_z-\lambda(r_d)$ annihilates each unwanted exponent and
leaves a nonzero multiple of the remaining one. Therefore the
coefficient of $e^{(d+1)z}$ must vanish separately. Since the two
relevant amplitudes have at most simple poles, their residues sum to
zero. This proves \eqref{eq:residue}.

Finally, use $q=t^{-1}$ and
$q^d-(d+1)=t^{-d}(1-(d+1)t^d)$ in \eqref{eq:mainW}.
The finite limiting value is exactly \eqref{eq:Lpositive}.
\end{proof}

\begin{theorem}[Exact exponential growth]\label{thm:radius}
For every $d\ge1$, the radius of convergence of $W_d(t)$ is exactly
$\rho_d=(d+1)^{-1/d}$, and
\begin{equation}\label{eq:limsup}
 \limsup_{k\to\infty} a_{d,k}^{1/k}=(d+1)^{1/d}.
\end{equation}
\end{theorem}
\begin{proof}
By \cref{thm:denominator,lem:separation}, there are no poles in
$|t|<\rho_d$. By \cref{lem:noncancel}, $t=\rho_d$ is an actual pole.
The radius is therefore exactly $\rho_d$.
The coefficient formula for the radius of convergence gives
\eqref{eq:limsup}.
\end{proof}

\subsection{The periodic leading term}
The same argument provides more than a limsup. Let
$\zeta_j=e^{2\pi i j/d}$ and $u_j=\rho_d\zeta_j$.
The function $L_d$ in \eqref{eq:L} is regular at all $u_j$:
\cref{lem:separation} leaves at most the single simple factor
$1-(d+1)t^d$ on this circle.

\begin{corollary}[Residue-class asymptotics]\label{cor:periodic}
There is a number $R>\rho_d$ such that
\begin{equation}\label{eq:periodicasymptotic}
 a_{d,k}=\rho_d^{-k}\mathcal A_{d,k\bmod d}+O(R^{-k}),
\end{equation}
where
\begin{equation}\label{eq:amplitudes}
 \mathcal A_{d,r}=\frac1d\sum_{j=0}^{d-1}L_d(u_j)\zeta_j^{-r}
 \quad(0\le r<d).
\end{equation}
Every $\mathcal A_{d,r}$ is real and nonnegative, and
$\sum_r\mathcal A_{d,r}=L_d(\rho_d)>0$.
\end{corollary}
\begin{proof}
At $t=u_j$, the contribution of the simple pole is
\[
 \frac{L_d(u_j)}{d}\,\frac{1}{1-t/u_j}.
\]
Subtracting these contributions removes all singularities on the first
pole circle. Since the function is rational, the remainder is analytic
in a strictly larger disk. A coefficient bound on an intermediate
circle gives the error term in \eqref{eq:periodicasymptotic}.
Extracting the simple-pole coefficients yields
\eqref{eq:amplitudes}. Along each residue class,
$\rho_d^k a_{d,k}$ tends to $\mathcal A_{d,r}$; nonnegativity of the
coefficients makes each limit nonnegative. Summing the discrete Fourier
formula over $r$ leaves only the term $j=0$.
\end{proof}

The theorem deliberately states a limsup for general $d$. Strict
positivity of every residue-class amplitude has not been proved here.
Nor is convergence of the consecutive-term ratio asserted; the explicit
$d=2$ formula shows that such a ratio limit can fail.

\section{The OEIS partition connection: a structural contrast}\label{sec:partitions}
Let $T(n,d)$ denote the A256193 partition triangle \cite{OEIS}.
The convention matters: after writing the parts of an ordinary partition
in nonincreasing order, choose exactly $d$ \emph{positions} to mark.
Different markings of equal-part positions are distinct. Equivalently,
\begin{equation}\label{eq:Tdefinition}
 T(n,d)=\sum_{\lambda\vdash n}\binom{\ell(\lambda)}{d},
\end{equation}
where $\ell(\lambda)$ is the number of parts. Its two-variable generating
function is
\begin{equation}\label{eq:Tgf}
 \sum_{n,d\ge0}T(n,d)x^ny^d
       =\prod_{m\ge1}\frac{1}{1-(1+y)x^m}.
\end{equation}
Indeed, a block of $r$ equal parts of size $m$ has marking polynomial
$(1+y)^r$, and summing over $r$ gives the $m$th geometric factor.

The initial diagonal agreement with $W_d$ does not extend to the
complete series. For example,
\begin{center}
\begin{tabular}{r|rrrrrrrr}
\toprule
$k$ & 0&1&2&3&4&5&6&7\\
\midrule
$a_{2,k}$ &1&4&11&31&65&157&298&661\\
$T(k+2,2)$&1&4&11&24&49&89&158&262\\
\bottomrule
\end{tabular}
\end{center}
The following result explains a global analytic distinction, without
requiring a combinatorial interpretation of each discrepancy.

\begin{theorem}[The partition diagonals are nonrational]\label{thm:partitioncontrast}
For every fixed $d\ge1$, put
\begin{equation}\label{eq:V}
 V_d(t)=\sum_{k\ge0}T(d+k,d)t^k.
\end{equation}
Then $V_d$ has radius of convergence~$1$ and is not rational. Moreover,
\begin{equation}\label{eq:growthcontrast}
 \limsup_{k\to\infty}
 \left(\frac{a_{d,k}}{T(d+k,d)}\right)^{1/k}=(d+1)^{1/d}>1.
\end{equation}
\end{theorem}
\begin{proof}
There are elementary bounds
\begin{equation}\label{eq:partitionbounds}
 p(k)\le T(d+k,d)\le(d+k)^d p(d+k).
\end{equation}
For the lower bound, append $d$ marked singleton parts to any unmarked
partition of $k$, placing them after the unmarked singleton parts.
This is injective. For the upper bound, each partition of $d+k$ has
at most $d+k$ positions available for marking.

For completeness, $p(n)$ is subexponential without needing a sharp
partition asymptotic. For any fixed $0<x<1$,
\[
 \sum_{n\ge0}p(n)x^n=\prod_{m\ge1}(1-x^m)^{-1}<\infty,
\]
since
\[
 \sum_{m\ge1}-\log(1-x^m)
 \le\sum_{m\ge1}\frac{x^m}{1-x^m}
 \le\frac{x}{(1-x)^2}.
\]
Thus $\limsup p(n)^{1/n}\le x^{-1}$; letting $x\uparrow1$ and using
$p(n)\ge1$ gives $p(n)^{1/n}\to1$. By \eqref{eq:partitionbounds},
$T(d+k,d)^{1/k}\to1$, proving the radius assertion and, together with
\cref{thm:radius}, \eqref{eq:growthcontrast}.

It remains to prove nonrationality. Partition numbers cannot have a
polynomial upper bound. If $p(n)=O(n^B)$, choose an integer $s>B+1$.
Taking each multiplicity of parts $1,\ldots,s$ independently from
$0,\ldots,M$ produces $(M+1)^s$ different partitions, each of total size
at most $N=Ms(s+1)/2$. But the assumed bound would imply
\[
 (M+1)^s\le\sum_{n=0}^{N}p(n)=O(N^{B+1})=O(M^{B+1}),
\]
a contradiction as $M\to\infty$.
By \eqref{eq:partitionbounds}, the coefficients of $V_d$ likewise admit
no polynomial upper bound. A rational series of radius~$1$ does have
such a bound: partial fractions express its coefficients as a finite
sum of polynomial factors times exponential factors of modulus at
most~$1$. Therefore $V_d$ cannot be rational.
\end{proof}

In particular, $W_d-V_d$ is nonrational and is not a polynomial.
The difference cannot be repaired by changing finitely many coefficients.
The earlier near-diagonal agreement concerns a finite initial region
for fixed $d$; the full tails belong to different generating-function
classes and have different exponential scales.
This conclusion does not claim an eventual coefficientwise inequality
for every $k$: \eqref{eq:growthcontrast} is the precise statement proved.

\section{Algorithms and independent exact verification}\label{sec:computation}
\subsection{Constructing the rational functions}
The program \texttt{code/derive.py} implements
\eqref{eq:g}--\eqref{eq:solveode}. A spectrum is a dictionary whose key is
the coefficient vector $(c_0,c_1,\ldots)$ of a frequency and whose value
is an exact rational function in $q$. Multiplication adds keys and
multiplies amplitudes; equal frequencies are combined before solving
the next differential equation.

At each level the program verifies, in $\Q(q)$, both the differential
identity and the initial condition:
\[
 (\lambda-q^d)c_{d,\lambda}=h_{d,\lambda},\qquad
 \sum_\lambda c_{d,\lambda}=0.
\]
It extracts $C_d$, reverses its numerator and denominator directly,
checks coprimality, and verifies that the resulting reduced denominator
divides the universal denominator. Direct coefficient reversal avoids
the expression swell of substituting $q=1/t$ into a large expanded
rational expression.

The finite frequency count $M_d$ is not by itself a bit-complexity bound.
Rational-function numerator and denominator sizes also grow, and
cancellation can be expensive. No polynomial-time claim in $d$ is made.
For a fixed derived denominator of degree $L$, however, successive tail
coefficients take $O(L)$ integer arithmetic operations each. Producing
$K$ coefficients takes $O(KL)$ such operations, apart from integer bit
costs. Standard polynomial reduction for a linear recurrence also permits
an isolated distant term to be computed in $O(L^2\log K)$ arithmetic
operations using naive polynomial multiplication.

\subsection{An independent path from the original recurrence}
The verifier \texttt{code/verify.py} uses only Python's standard library.
It does not call SymPy and does not reconstruct the finite spectra.
Instead it evaluates the original recurrence after reversal, truncating
in deficit degree. The degree gaps in \cref{sec:stabilization} show that
all shifts are nonnegative, so discarding higher deficit coefficients
cannot affect lower ones.

For every $1\le d\le6$ and $0\le k\le40$, it compares the rational-series
coefficient against the original recurrence at three lengths:
\[
 n=d+k+1,\qquad n=d+k+2,\qquad n=48.
\]
These are $6\cdot41\cdot3=738$ exact integer comparisons.
The first length is the proved stabilization boundary; the other two
check that the same coefficient persists. In addition, all 251 exported
coefficients of $W_1$ and $W_2$ are checked against their independent
closed formulas. No floating-point arithmetic is involved.

\begin{table}[htbp]
\centering\small
\caption{A prefix recomputed by both methods. Entries are $a_{d,k}$.}
\label{tab:prefix}
\begin{tabular}{r|rrrrrr}
\toprule
$k$ & $d=1$ & $d=2$ & $d=3$ & $d=4$ & $d=5$ & $d=6$\\
\midrule
0 &1&1&1&1&1&1\\
1 &3&4&5&6&7&8\\
2 &7&11&16&22&29&37\\
3 &15&31&41&63&92&129\\
4 &31&65&112&155&247&376\\
5 &63&157&244&393&590&966\\
6 &127&298&542&869&1404&2258\\
7 &255&661&1155&1863&3079&5140\\
8 &511&1196&2288&3886&6471&11078\\
9 &1023&2516&4552&7863&13216&22921\\
10&2047&4434&8852&15217&26354&46011\\
11&4095&9002&16403&29134&50977&90269\\
\bottomrule
\end{tabular}
\end{table}

The symbolic formulas prove statements for all indices. The finite
comparisons serve a different purpose: they guard against transcription,
normalization, implementation, and indexing errors. The computations are
not used as a substitute for the general arguments.

\subsection{Reproduction and files}
From the archive root, the following commands reproduce the exact results:
\begin{lstlisting}
python -m pip install -r requirements.txt
python code/derive.py --max-d 6 --check-only
python code/verify.py --max-k 40 --export-count 251
pdflatex -interaction=nonstopmode -halt-on-error article.tex
pdflatex -interaction=nonstopmode -halt-on-error article.tex
\end{lstlisting}
The derivation was run with SymPy 1.14.0. The independent verifier needs
Python 3.9 or later and no third-party package. The data directory contains
normalized numerator and denominator arrays in ascending degree order,
exact spectral certificates, CSV tables, individual coefficient files,
and the actual verification logs. No Lean or other proof-assistant
formalization is claimed.

\section{What is settled, and what remains}\label{sec:conclusion}
The whole-series rationality problem has a uniform affirmative answer,
with explicit finite algebraic data. The proof depends only on the exact
recurrence \eqref{eq:original}. It establishes an integer recurrence for
every descent count, an exact radius of convergence, and periodic
residue-class asymptotics. The two-descent case is completely explicit
in elementary sequences.

Several sharper questions remain beyond this report. A general rule for
the cancellations in $\mathcal D_d$ would determine minimal recurrence
orders without symbolic gcd computations. A proof that every
$\mathcal A_{d,r}$ is strictly positive would upgrade the general limsup
to an ordinary $k$th-root limit. A direct combinatorial explanation of the
rational numerators or the complete discrepancy $a_{d,k}-T(d+k,d)$ could
make the algebra more transparent. Finally, the growth and structure of
the bivariate series $\sum_{d\ge1}u^dW_d(t)$ is a separate problem; fixed-$d$
rationality does not imply its rationality in two variables.

The useful outcome of this investigation is therefore not a fitted
recurrence or a large computational prefix. It is a finite-spectrum
mechanism that determines each complete stabilized generating function,
together with a precise distinction between those series and the OEIS
partition diagonals that originally suggested them.

\appendix
\section{The spectral construction in pseudocode}\label{app:algorithm}
In the following pseudocode, dictionary keys are polynomials in $q$ and
values lie in $\Q(q)$. Zero entries are removed after every collection.
The single-symbol key $q^d$ refers to a frequency, not to its amplitude.
\begin{lstlisting}
F[0] := { 1 : 1 }
for d := 1, ..., D:
    G := empty dictionary
    for a := 0, ..., d-1:
        b := d-1-a
        for (lambda, c) in F[a]:
            for (mu, e) in F[b]:
                G[lambda+mu] += q^b * c * e
    for (lambda, c) in F[d-1]:
        G[lambda] -= q^(d-1) * c

    F[d] := empty dictionary
    for (lambda, h) in G:
        F[d][lambda] := h / (lambda-q^d)
    F[d][q^d] := -sum(values of F[d])

    C[d](q) := F[d][q^d]
    W[d](t) := t^(-d*(d+1)) * C[d](1/t)
    cancel and normalize the denominator to have constant term 1
\end{lstlisting}
The decisive invariant is $\deg\lambda<d$ for every key in $G$.
It guarantees that the division is legal over $\Q(q)$ and that the
$q^d$ mode is the unique mode surviving the formal leading-tail limit.

\section{A compact audit of claims}\label{app:audit}
\begin{longtable}{@{}p{0.27\textwidth}p{0.65\textwidth}@{}}
\toprule
Claim & Basis and qualification\\
\midrule
\endhead
Definition of the polynomial family & The published recurrence
\cite[Theorem 4.1]{ACS}, reproduced as \eqref{eq:original}.\\[4pt]
Existence of $W_d$ & Known stabilization result \cite{ACS}; independently
proved from the recurrence in \cref{prop:stabilization}.\\[4pt]
Rationality for every $d$ & General proof in
\cref{thm:spectrum,thm:rationality,thm:denominator}; not inferred from
experiments.\\[4pt]
Exact growth radius & Pole exclusion and a noncancellation proof in
\cref{lem:separation,lem:noncancel,thm:radius}.\\[4pt]
Closed formulas for $d=2$ & Exact differential calculation, rational
identity, and coefficient extraction in \cref{sec:examples}.\\[4pt]
Denominator degrees through $d=6$ & Exact rational-function arithmetic,
coprimality checks, and retained machine-readable results.\\[4pt]
Coefficient verification & 738 comparisons with the original recurrence,
plus closed-form checks; all passed. This is an implementation audit,
not proof-assistant verification.\\[4pt]
Novelty & No all-$d$ rationality theorem was found in the three primary
papers reviewed. No claim is made that the entire literature has been
exhaustively checked.\\
\bottomrule
\end{longtable}

\begin{thebibliography}{9}
\bibitem{DGGs}
W. Dugan, S. Glennon, P. E. Gunnells, and E. Steingr\'imsson,
\emph{Tiered trees, weights, and $q$-Eulerian numbers},
Journal of Combinatorial Theory, Series A \textbf{164} (2019), 24--49.
\href{https://arxiv.org/abs/1702.02446}{arXiv:1702.02446}.
The motivating discussion of stabilized series and partition diagonals
is Section~6.12.

\bibitem{ACS}
A. Agrawal, C. Choi, and N. Sun,
\emph{On Permutation Weights and $q$-Eulerian Polynomials},
Annals of Combinatorics \textbf{24} (2020), 363--378.
\href{https://doi.org/10.1007/s00026-020-00493-5}{doi:10.1007/s00026-020-00493-5}.
\href{https://arxiv.org/abs/1809.07398}{arXiv:1809.07398}, version~4.

\bibitem{HT}
E. Huang and R. Tang,
\emph{Minimum Decomposition on Maxmin Trees},
\href{https://arxiv.org/abs/2310.14385}{arXiv:2310.14385} (2023), version~2.
The end of Section~5 discusses the challenge of reaching coefficients
beyond the near-diagonal partition correspondence.

\bibitem{OEIS}
The OEIS Foundation Inc.,
\href{https://oeis.org/A256193}{A256193: Number $T(n,k)$ of partitions of $n$
into two sorts of parts having exactly $k$ parts of the second sort},
The On-Line Encyclopedia of Integer Sequences.
The marking convention is made explicit in \eqref{eq:Tdefinition}.
\end{thebibliography}
\noindent\small\color{Muted}Primary sources and the OEIS entry were consulted on
September 19, 2026. All displayed new formulas and computational tables
were derived or recomputed for this report.
\end{document}
